\section{The finite Adams--Novikov calculation}\label{sec:finite}

This section proves the Adams--Novikov assertions used in
Section~\ref{sec:cube}.  The calculation determines only the kernels,
images, and products needed in stem $322$; it does not require a complete
description of either adjacent stable stem.  All vector spaces in this
section are over $\F_3$.

\subsection{The product and its algebraic conventions}

Let $M=S/3$, put $B=\BP_*/3$, and write
\[
 D_{f,d}=\Ext^{f,4d}_{(B,\BP_*\BP/3)}(B,B)=E_2^{f,4d}(M).
\]
Thus $f$ is Adams--Novikov filtration, $d$ is one quarter of internal degree, and
the corresponding stable stem is $4d-f$.  These groups vanish unless the
internal degree is divisible by four.  Consequently
\[
 E_5=D,\qquad E_6=E_9=H(D,d_5),\qquad
 |d_5|=J=(5,1),\quad |d_9|=R=(9,2).
\]
When no confusion is possible, we abbreviate $H(D,d_5)$ to $H$.
The signs in a Leibniz rule are $(-1)^f$, since $4d-f$ has parity $f$.

We first give the tower comparison needed for the multiplicative
spectral sequence.  It also identifies the totalization convention
with the Adams stages used in Section~\ref{sec:cube}.

\begin{lemma}[Unit cubes and the pairing]\label{finite:pairing}
Let $I=\operatorname{fib}(S\longrightarrow\BP)$ and
$F_qZ=I^{\wedge q}\wedge Z$.  For the augmented Amitsur cosimplicial
spectrum
\[
 X^n(Z)=\BP^{\wedge(n+1)}\wedge Z
\]
there are natural equivalences
\[
 \operatorname{fib}\bigl(Z\longrightarrow\Tot_m X^\bullet(Z)\bigr)
 \simeq F_{m+1}Z,
\]
compatible with the maps as $m$ varies.  A unital binary multiplication
$\mu\colon M\wedge M\longrightarrow M$ induces pairings on its
Adams--Novikov spectral sequence with the Leibniz rule on every page.
The pairing on $E_2(M)$ is the ordinary graded-commutative Ext product.
These statements do not require an associative multiplication on $M$.
\end{lemma}

\begin{proof}
Let $\mathcal P_m$ be the poset of nonempty subsets of
$[m]=\{0,\ldots,m\}$.  The functor
\[
 u\colon\mathcal P_m\longrightarrow\Delta_{\leq m},\qquad
 S\longmapsto[|S|-1]
\]
sends an inclusion to the induced order-preserving injection.  We show
that $u$ is homotopy initial.  For $0\leq k\leq m$, the category
$(u\downarrow[k])$ is the face poset of the simplicial complex
$K_{m,k}$ whose nonempty faces are weakly increasing partial maps
$[m]\dashrightarrow[k]$.  Indeed, a face consists of a subset $S$ and
an order-preserving map from its induced ordinal to $[k]$.

The complex $K_{m,k}$ is contractible when $k\leq m$.  To see this,
start with the subcomplex whose last argument is either undefined or
has value $k$.  This is a cone on $K_{m-1,k}$, with cone vertex $(m,k)$.
For each $a<k$, the remaining faces with last value $a$ form the cone
with vertex $(m,a)$ on $K_{m-1,a}$.  Its intersection with the previously
constructed complex is exactly that base.  Since $a\leq m-1$, the base
is contractible by induction.  Attaching a cone on a contractible
subcomplex preserves contractibility.  The initial case $m=k=0$ is a
point.  This proves the assertion, including the case $k=m$, where the
first cone is contractible without any assumption on its base.
The homotopy-initiality criterion therefore identifies $\Tot_m X$ with
the homotopy limit of the nonempty unit cube indexed by $\mathcal P_m$.
Including its empty vertex gives the cube with initial value $Z$ and
one unit insertion in each of its $m+1$ directions.  Its total fibre is
$I^{\wedge(m+1)}\wedge Z$, because smash products of spectra preserve
finite fibres in each variable.  This construction is compatible with
restriction of cubes and hence with the tower maps.

Choose the $E_4$ structure on $\BP$ furnished by
\cite[Theorem~1.1]{BasterraMandell}.  In particular, its associative algebra
structure in associative algebras supplies coherent pointwise pairings
of the Amitsur construction.  The external pairings, followed by
$\mu$, give the pairing of totalization spectral sequences, and thus
of the identified Adams towers.  The boundary formula for a pairing
of filtered objects is precisely the Leibniz rule; it uses the binary
pairing and its compatibility with the tower, not associativity of
$\mu$.  The required unital $\mu$ is provided by
Lemma~\ref{cube:endomorphisms}.

Finally, the sphere-input Amitsur levels have even, $3$-torsion-free
homotopy groups.  Smashing with $M$ reduces these groups modulo $3$.
Unitality makes the resulting coefficient pairing
$B\otimes_{\BP_*}B\longrightarrow B$ ordinary multiplication.
Since $3$ is invariant in the Hopf algebroid, the normalized cobar
complex after this reduction computes $D$.  Its induced cup product
is the ordinary Ext product, proving the last assertion.  In
particular, associativity and graded commutativity hold on $E_2$ and
pass to all subsequent pages by taking homology.
\end{proof}

We use Hazewinkel generators $v_i$ and coordinates $t_i$, with
$|v_i|=|t_i|=2(3^i-1)$ in internal degree.  Thus their weight is
$(3^i-1)/2$.  The logarithm convention is
\[
 \eta_R(l_i)=\sum_{j=0}^i l_jt_{i-j}^{3^j},\qquad l_0=t_0=1.
\]
The integral recursions and their reductions are given in
Appendix~\ref{app:recurrence}.  In characteristic three they give
\begin{equation}\label{finite:right-unit}
 \eta_R(v_1)=v_1,\qquad
 \eta_R(v_2)=v_2+v_1(t_1^3-v_1^2t_1).
\end{equation}
The primitive coefficient ring is exactly $\F_3[v_1]$.  For
completeness, suppose a primitive polynomial depends on $v_s$, with
$s\geq2$ maximal.  Set all positive-index $t_i$ except $t_{s-1}$ to
zero.  The recursion fixes $v_1,\ldots,v_{s-1}$ and sends
\[
 v_s\longmapsto v_s+v_1t_{s-1}^3-v_1^{3^{s-1}}t_{s-1}.
\]
If the degree of the polynomial in $v_s$ is $N>0$, its transformed
expression has a nonzero coefficient on $t_{s-1}^{3N}$: this
coefficient is its leading coefficient in $v_s$, multiplied by
$v_1^N$.  The coefficient ring is a domain.  This contradicts
primitivity and proves the assertion by descent on $s$.

Put $v=v_1$ and $a=[t_1]$.  For a coefficient $q$ whose right-unit
difference is divisible by $v_1^m$ modulo three, write
\[
 \delta_m(q)=\frac{\eta_R(q)-q}{v_1^m}\pmod3.
\]
Reduction precedes division.  Equation~\eqref{finite:right-unit} and
the invariance of $v_1$ show that the following expressions are
one-cocycles:
\begin{align*}
 h&=\delta_1(v_2),& r&=\delta_3(v_2^3),&
 s'&=\delta_1(v_2^2),\\
y&=\delta_9(v_2^9-v_1^8v_2^7),&&&
 t_i'&=\delta_1(v_2^i).
\end{align*}
Their bidegrees are respectively
\begin{gather*}
 |v|=(0,1),\quad |a|=(1,1),\quad |h|=(1,3),\quad |r|=(1,9),\\
 |s'|=(1,7),\quad |y|=(1,27),\quad |t_i'|=(1,4i-1).
\end{gather*}
For example, if $w=t_1^3-v_1^2t_1$, then
$(v_2+v_1w)^9-v_2^9=v_1^9w^9$; the difference of the seventh powers
is divisible by $v_1$.  This proves the divisibility defining $y$ and
also the useful identity
\begin{equation}\label{finite:corrected-y}
 \delta_9(v_2^9)=y+t_7'.
\end{equation}

The notation $B(x)$ denotes a Bockstein; the unadorned $B$ still
denotes the coefficient ring.  For a one-cocycle $x$ define its
algebraic Bockstein class by
\begin{equation}\label{finite:bockstein}
 B(x)=\frac{\Delta\widetilde x-\widetilde x\otimes1
                  -1\otimes\widetilde x}{3}\pmod3,
\end{equation}
where $\widetilde x$ is an integral lift.  Coefficients are written on
the left: $\Delta$ is left-linear, and moving coefficients in the last
term into this convention uses $\eta_R$.  The numerator is divisible
by $3$ because $x$ is a cocycle.  Changing the lift changes the result
by a coboundary.  With the usual normalized-cobar differential,
$B$ is the negative of the reduced connecting homomorphism for
multiplication by $3$.  Set
\[
 b=B(h),\qquad c'=B(r),\qquad P=aB(t_4'),
 \qquad |b|=(2,3),\ |c'|=(2,9),\ |P|=(3,16).
\]
These classes are, up to units, the reductions of the sphere classes
$\beta_1$, $\beta_{3/3}$, and $\alpha_1\beta_4$.  The beta
coboundaries defining them come respectively from $v_2$ modulo
$(3,v_1)$, $v_2^3$ modulo $(3,v_1^3)$, and $v_2^4$ modulo $(3,v_1)$.
The classical Adams--Novikov differentials and naturality therefore
give
\begin{equation}\label{finite:classical-differentials}
 d_5(c')=\epsilon ab^3,\qquad P\in E_9(M),\qquad
 d_9(P)=\epsilon' b^6,
 \qquad \epsilon,\epsilon'\in\F_3^\times.
\end{equation}
We use these classical results in the form recalled in
\cite[Lemma 2.5]{BelmontShimomura}; see also \cite{Ravenel}.
Their natural images in the Moore spectral sequence are allowed to
be zero or boundaries.  No assertion that such an image is nonzero
is needed below.  The classes $a$ and $b$ come from the permanent
classes $\alpha_1$ and $\beta_1$.

Here is the conclusion of the section.  The coefficient in its
hypothesis means the coefficient in a normalized one-cobar
representative after setting all positive-degree coefficient
generators $v_i$ to zero.

\begin{lemma}[Finite Adams--Novikov lemma]\label{finite:main}
Suppose $z\in E_2^{1,108}(M)=D_{1,27}$ is a permanent cycle whose
coefficient on $t_1^{27}$ is $1$.  Then:
\begin{enumerate}[label=\textup{(\roman*)}]
 \item $z\equiv y\pmod{\F_3v^{26}a}$ and $B(y)^3=0$ in $D$;
 \item multiplication by $v$ is zero from
 $E_6^{10,328}(M)$ to $E_6^{10,332}(M)$.  The same holds on
 $E_9$ and $E_{10}$;
 \item $E_{10}^{f,322+f}(M)=0$ for every
 $f\equiv2\pmod4$ with $f\geq14$.
\end{enumerate}
\end{lemma}

\subsection{The finite affine family of possible differentials}

The three rank calculations in Table~\ref{finite:low-ranks} fix the
low-dimensional classes and the cubic relation.  In the cochain complex of
Appendix~\ref{app:recurrence}, let $n_-,n,n_+$ be the dimensions of
the three terms surrounding the indicated bidegree.  Let $r_-$ and
$r$ be the incoming and outgoing coboundary ranks.  For a displayed
list $C$ of cocycle representatives, let $r_C$ be the rank obtained
by adjoining those columns to the incoming coboundary matrix.

\begin{table}[htbp]
\centering
\(
\begin{array}{cc|rrr|rr|r|l}
 f&d&n_-&n&n_+&r_-&r&r_C&C\\\hline
 1&27&12&49&94&11&35&14&v^{26}a,\ y,\ t_7'\\
 6&10&3&2&1&1&0&2&vb^3\\
 6&81&9301&11419&12864&5256&6156&5256&B(y)^3
\end{array}
\)
\caption{Cochain ranks identifying the distinguished classes.}
\label{finite:low-ranks}
\end{table}

The construction and correctness of these matrices are proved in
Proposition~\ref{alg:resolution-correct}; their finite reduction, and
those used below, are specified in
Proposition~\ref{alg:table-verification}.  In particular, the table is
an exact cochain calculation over $\F_3$, rather than a choice of
names for classes in a chart.  Since cohomology has dimension
$n-r_--r$, it proves
\begin{equation}\label{finite:low-classes}
 D_{1,27}=\langle v^{26}a,y,t_7'\rangle,\qquad
 D_{6,10}=\langle vb^3\rangle\cong\F_3,\qquad B(y)^3=0.
\end{equation}
The three basis elements in $D_{1,27}$ have $t_1^{27}$ coefficients
$0,1,0$.  The coefficient is unchanged by a one-cobar boundary:
specializing the positive coefficient generators to zero makes every
right-unit difference zero.  Thus the hypothesized $z$ has the form
\begin{equation}\label{finite:z-form}
 z=y+e\,t_7'+\mu v^{26}a,\qquad e,\mu\in\F_3.
\end{equation}

Write $L=\epsilon^{-1}d_5$, in product coordinates normalized by
\eqref{finite:classical-differentials}.  It is a square-zero
derivation of bidegree $J$.  The elementary vanishing bound
\begin{equation}\label{finite:elementary-vanishing}
 D_{f,d}=0\quad\text{if}\quad d<\frac32f-\frac12
\end{equation}
follows from the additive-digit cobar filtration proved in
Lemma~\ref{finite:highweight} below: the only generator with positive
deficit $\frac32f-d$ is the exterior class of weight one, with
deficit $1/2$.  Applying this bound to the targets, together with
permanence of $a$ and $b$, gives
\[
 L(v)=L(a)=L(b)=L(h)=L(s')=0,\qquad L(c')=ab^3.
\]
Also $D_{0,*}=\F_3[v]$, so $L$ vanishes in filtration zero.
In fact $v$ is permanent: each possible outgoing differential has
target $(r,1+(r-1)/4)$, which vanishes by
\eqref{finite:elementary-vanishing} for $r\geq2$, and filtration
zero admits no incoming differential.
Equation~\eqref{finite:low-classes} implies
\[
 L(r)=k\,vb^3\quad(k\in\F_3).
\]
Since $z$ is permanent and $L(v^{26}a)=0$, Equation
\eqref{finite:z-form} imposes $L(y+e\,t_7')=0$.

We now encode only necessary conditions on this actual differential.
For a cutoff $N$, use the finite grid
\begin{equation}\label{finite:grid}
 \mathcal G_N=\{(f,d):0\leq d\leq N,\ 0\leq f\leq32,
                                  \ 2d-3f\geq0\}.
\end{equation}
Put independent unknowns in every entry of
$L_T\colon D_T\longrightarrow D_{T+J}$ for which both slots belong
to the grid.  Impose its vanishing in filtration zero.  For each
\[
 x\in\{a,v,h,b,c',r,s',y+e\,t_7'\},
\]
impose, for every $q\in D_T$, the Leibniz identity
\begin{equation}\label{finite:leibniz-equations}
 L(xq)=L(x)q+(-1)^{f_x}xL(q)
\end{equation}
whenever the four slots $T,T+J,T+|x|,T+|x|+J$ are present.  The
known values of $L(x)$ enter as constants.  For $s'$ and
$y+e\,t_7'$, retain only multiplication maps whose output has
weight at most $91$.  This restriction deliberately permits unused
slots and equations; every equation retained still holds for the
actual $L$.  In particular, classes on the boundary line excluded
by \eqref{finite:grid} are not assigned additional unknown maps.

At cutoff $60$, these are affine linear equations for each of the
nine choices $(k,e)$.  Their coefficient rank is $2667$ in every
case.  The augmented rank is $2667$ for $(k,e)=(1,0)$ and $2668$
for each of the other eight choices.  Hence
\begin{equation}\label{finite:forced-values}
 k=1,\qquad e=0.
\end{equation}
This proves the first assertion about $z$ in
Lemma~\ref{finite:main}; the cubic assertion was already proved in
\eqref{finite:low-classes}.  Notice that the correction in
\eqref{finite:corrected-y} is essential to this statement.

At cutoff $92$ keep $(k,e)=(1,0)$ and also impose the necessary
\emph{linear consequences} of $L^2=0$.  More precisely, substitute
the current affine parametrization in each available equation
$L_{T+J}L_T=0$.  Row-reduce the resulting polynomial equations with
all quadratic monomials ordered before the linear monomials and
the constant.  Rows in which every quadratic term vanishes give
affine equations satisfied by every square-zero point.  Adjoin
these rows to the linear system, parametrize again, and repeat until
no new linear condition appears.  The resulting ranks are listed in
Table~\ref{finite:affine-ranks}.

\begin{table}[htbp]
\centering\small
\(
\begin{array}{c|r|r|c|r}
 N&\text{unknown entries}&\text{initial affine rank}&
     \text{new affine constraints}&
     \text{final affine dimension}\\\hline
 60&2779&2667&\text{not imposed}&112\\
 92&23384&23119&38,\ 1,\ 0&226
\end{array}
\)
\caption{Necessary affine conditions on the normalized fifth differential.}
\label{finite:affine-ranks}
\end{table}

The cutoff-$92$ affine family may contain nonsquare-zero maps.  What
has been proved, and all that is needed, is that it contains the
restriction of the actual $d_5$ after normalization.  Neither
uniqueness of $d_5$ nor a choice among the remaining $226$ parameters
is assumed.

\begin{remark}[Product signs]\label{finite:signs}
The products used to obtain these matrices are represented by signed
chain lifts $\phi$ with
$\partial\phi=(-1)^f\phi\partial$ for a class of filtration $f$;
the fixed class is lifted first.  As explained in
Appendix~\ref{app:recurrence}, comparison with cobar cup products
can use the opposite ordering of the factors, giving the sign twist
$(-1)^{fg}$ in filtrations $f,g$.  The two conventions are identified by
rescaling filtration $f$ by $(-1)^{\binom f2}$.  Transporting the
differential through this identification gives the corresponding
row signs, and a further unit normalizes $L(c')=ab^3$.  All displayed
systems use this consistent convention.  Such changes preserve the
kernels, images, ranks, and lines of possible $d_9$ values used
below.
\end{remark}

\subsection{Universal bounds for kernels and images}

For a projected matrix of the final affine family, write
$A(u)=A_0+\sum_{i=1}^r u_iA_i$.  Parameters may be correlated across
different matrices.  Define
\begin{align}
 K^-&=\bigcap_{i=0}^r\ker A_i,&
 K^+&=A_0^{-1}\left(\sum_{i=1}^r\im A_i\right),\label{finite:kernel-bounds}\\
 I^-&=A_0\left(\bigcap_{i=1}^r\ker A_i\right),&
 I^+&=\sum_{i=0}^r\im A_i.\label{finite:image-bounds}
\end{align}
For an empty set of parameter matrices, the intersection is the
whole source and the sum is zero.

\begin{lemma}\label{finite:bounds-valid}
For every parameter value $u$ one has
\[
 K^-\subseteq\ker A(u)\subseteq K^+,
 \qquad I^-\subseteq\im A(u)\subseteq I^+.
\]
\end{lemma}
\begin{proof}
Every $A_i$ annihilates a vector in $K^-$.  If $A(u)x=0$, then
$A_0x=-\sum_{i>0}u_iA_ix$, proving the second kernel inclusion.
If all $A_i$ for $i>0$ annihilate $x$, then $A(u)x=A_0x$;
this proves the first image inclusion.  The last follows directly
from the formula for $A(u)$.  No independence assumption on the
parameters is used.
\end{proof}

At a slot $p$, let $K_p^\pm$ denote these bounds for the map out of
$p$, and $I_p^\pm$ the bounds for the map into $p$.  Negative
filtration is zero.  If an otherwise relevant pair is unavailable,
use zero as a lower bound and the whole source or target as the
upper bound.  Table~\ref{finite:enclosures} records the resulting
dimensions; its column $D$ means $\dim D_{f,d}$.

\begin{table}[htbp]
\centering\small
\(
\begin{array}{cc|r|rrrr@{\qquad}cc|r|rrrr}
 f&d&D&I^-&I^+&K^-&K^+&f&d&D&I^-&I^+&K^-&K^+\\\hline
 4&69&7&0&0&4&4&13&71&7&3&3&3&4\\
 12&76&9&4&5&6&6&21&78&4&2&2&2&2\\
 5&82&5&0&0&4&4&10&82&8&1&5&6&6\\
 9&83&6&1&1&2&2&10&83&6&1&1&2&4\\
 14&84&7&4&4&5&5&7&85&11&0&2&8&9\\
 18&85&8&3&4&5&5&22&86&4&2&2&2&2\\
 23&86&7&3&3&3&4&16&87&5&2&2&3&3\\
 27&87&4&2&2&2&2&15&92&6&2&6&0&6
\end{array}
\)
\caption{Universal lower and upper bounds for images and kernels.}
\label{finite:enclosures}
\end{table}

The two additional span calculations are
\begin{equation}\label{finite:span-dimensions}
 \dim\bigl(vK^+_{10,82}+I^-_{10,83}\bigr)=1,
 \qquad
 \dim\bigl(K^+_{22,86}+I^-_{22,86}\bigr)=2.
\end{equation}
As the corresponding $I^-$ spaces have dimensions $1$ and $2$,
these equations give
\begin{equation}\label{finite:first-consequences}
 \bigl(v\colon H_{10,82}\longrightarrow H_{10,83}\bigr)=0,
 \qquad H_{22,86}=0.
\end{equation}
Indeed, $v$ sends every actual cycle in $K^+_{10,82}$ into
$I^-_{10,83}$, an actual boundary; at $(22,86)$ every actual cycle
is already in the indicated lower bound for boundaries.  The first
identity denotes the zero map, not a claim that $H_{10,83}$ itself
vanishes.  Since $v$ is permanent, this proves
Lemma~\ref{finite:main}\textup{(ii)} on all the stated pages.

\subsection{A uniform estimate in high filtration}

The next estimate supplies the two product maps that lie beyond
the finite affine calculation and handles the rest of the stem.

\begin{samepage}
\begin{lemma}[High-weight product bound]\label{finite:highweight}
For $k\geq1$, multiplication by $b^k$ into $D_{f,d}$ is surjective if
\begin{equation}\label{finite:surjective-bound}
 d<S_k(f):=\frac92f+\frac12-6k,
\end{equation}
and injective if
\begin{equation}\label{finite:injective-bound}
 d<I_k(f):=\frac92f-4-6k=S_k(f-1).
\end{equation}
\end{lemma}
\end{samepage}

\begin{proof}
Use the ordinary normalized cobar complex.  Order its monomials
first by increasing number of coefficient factors $v_i$, counted
with multiplicity; next by decreasing total coefficient weight;
and finally by the decreasing lexicographic order on the total
$t_i$ exponent tuple, with the largest index dominant.  Filter by
tails of this order.  At a fixed internal weight this filtration
is finite, including across all cobar lengths.

The ideal generated by the positive-degree coefficients is
invariant under the right unit.  Thus a right-unit correction cannot
decrease the coefficient count.  If it preserves that count and
introduces a nonconstant $t$ monomial, it decreases the coefficient
weight.  After setting the coefficients to zero, each nonadditive
coproduct term of $t_i$ uses lower-index $t$'s and hence decreases
the final tuple.  These assertions follow directly from the
Hazewinkel recursion of Appendix~\ref{app:recurrence}.  They show
that the associated graded cobar differential is the one for
ordinary additive coordinates with primitive coefficients.

The cohomology of this associated graded complex is
\begin{equation}\label{finite:additive-page}
 \F_3[v_1,v_2,\ldots]\otimes
 \bigotimes_{i\geq1,\,j\geq0}
 \bigl(\Lambda(e_{i,j})\otimes\F_3[c_{i,j}]\bigr),
\end{equation}
where, for $w_{i,j}=3^j(3^i-1)/2$,
\[
 |e_{i,j}|=(1,w_{i,j}),\qquad |c_{i,j}|=(2,3w_{i,j}).
\]
One obtains this formula by writing each additive coordinate in
base-three digits.  Its graded dual is a tensor product of
truncated polynomial algebras $\F_3[x]/(x^3)$, with $x$ of weight
$w_{i,j}$.  The alternating free resolution with maps $x$ and
$x^2$ has one exterior Ext generator in degree $(1,w_{i,j})$ and
one polynomial generator in degree $(2,3w_{i,j})$.  Tensoring these
resolutions gives \eqref{finite:additive-page}; only finitely many
factors occur in any fixed weight.

The leading class of $b$ on this page is $c_{1,0}$.  Indeed,
$h=t_1^3-v_1^2t_1$, and its integral Bockstein has coefficient-free
part
\[
 t_1^2\mid t_1+t_1\mid t_1^2,
\]
the standard polynomial generator for the weight-one digit.
All the other terms have positive coefficient count.

For a monomial on \eqref{finite:additive-page} define its deficit
to be $\frac92f-d$.  Among exterior generators the only positive
deficits occur at weights $1,3,4$, and they are respectively
$7/2,3/2,1/2$.  As these generators are exterior, their total
contribution is at most $11/2$.  The polynomial generator $b$
has deficit $6$; every other polynomial generator and every
coefficient factor has nonpositive deficit.  Consequently a
monomial not divisible by $b^k$ has deficit at most
\[
 \frac{11}{2}+6(k-1)=6k-\frac12.
\]
It follows that the quotient of \eqref{finite:additive-page} by
$b^k$ vanishes in bidegree $(f,d)$ when $d<S_k(f)$.
Multiplication by $b^k$ on that page is injective everywhere.

To transfer these assertions to $D$, take the filtered mapping cone
$Q$ of the cochain cup map
\[
 b^k\colon C^{*-2k,d-3k}\longrightarrow C^{*,d}.
\]
Shift the source filtration by the leading monomial of $b^k$.
The first cohomology page of the cone is exactly the quotient just
described, since its leading cup map is injective.  The filtration
is finite in fixed weight, so $H^{f,d}(Q)=0$ for $d<S_k(f)$.
The long exact sequence contains
\[
 H^{f-1,d}(Q)\longrightarrow D_{f-2k,d-3k}
 \xrightarrow{\ b^k\ }D_{f,d}\longrightarrow H^{f,d}(Q).
\]
Vanishing of the right term proves surjectivity.  Vanishing of the
left term uses $d<S_k(f-1)$ and proves injectivity, including the
one-filtration shift in \eqref{finite:injective-bound}.
\end{proof}

We spell out how ranks on $D$ pass to its $L$-homology.  Since $b$
has even filtration and $L(b)=0$, every power $q=b^k$ is a map of
the complexes defined by $L$.

\begin{lemma}[Lifting cycles and boundaries]\label{finite:homology-transfer}
Let $q=b^k$ and let $t$ be a target slot.
\begin{enumerate}[label=\textup{(\roman*)}]
 \item If multiplication by $q$ is surjective into $D_t$ and
 injective into $D_{t+J}$, it is surjective into $H_t$.
 \item If it is injective into $D_t$ and surjective into
 $D_{t-J}$, it is injective into $H_t$.
\end{enumerate}
\end{lemma}
\begin{proof}
For (i), lift a cycle $c\in D_t$ to $x$ with $qx=c$.
Then $qLx=Lc=0$, and injectivity into $D_{t+J}$ gives $Lx=0$.
For (ii), suppose $x$ is a cycle and $qx=Ly$ in $D_t$.
Surjectivity into $D_{t-J}$ gives $y=qz$.  Hence
$q(x-Lz)=0$; injectivity into $D_t$ gives $x=Lz$.
\end{proof}

\subsection{Vanishing after the ninth differential}

We now prove Lemma~\ref{finite:main}\textup{(iii)}.  First,
$H_{22,86}=0$ by \eqref{finite:first-consequences}, so its
$d_9$-homology vanishes as well.

The remaining low-filtration vanishings will use
$d_9(P)=\epsilon'b^6$ to turn cycles into boundaries.  We seek
positive integers $h_0,\ell$ with $h_0+\ell=6$ and classes $[x],[w]$
in $H$ such that a given $d_9$-cycle $c$ satisfies
\[
 c=b^{h_0}[x],\qquad P[x]=b^\ell[w].
\]
If multiplication by $b^{h_0}$ is injective from the slot of
$d_9[x]$, then $d_9[x]=0$.  Applying $d_9$ to the second
factorization gives $b^\ell d_9[w]=\epsilon'b^\ell c$.
A second injection, this time by $b^\ell$ on the slot containing
$c$, makes $c$ a boundary.  We first encode these factorizations
using bounds valid for the actual $L$; the required injection
tests follow.

Addition and subtraction of slots always mean bidegree
addition and subtraction; an element in a subscript denotes its
bidegree.  For $m=(14,84)$ and $(18,85)$ respectively set
\[
 h_0=5,3,\qquad \ell=6-h_0,\qquad
 u=m-h_0|b|,\quad p=u+|P|,\quad z=p-\ell|b|.
\]
Thus the complete degree data are
\[
\begin{array}{c|cc|ccc}
 m&h_0&\ell&u&p&z\\\hline
 (14,84)&5&1&(4,69)&(7,85)&(5,82)\\
 (18,85)&3&3&(12,76)&(15,92)&(9,83).
\end{array}
\]
For these choices define
\begin{equation}\label{finite:cancellation-spaces}
 A=\{x\in K^-_u:Px\in I^-_p+b^\ell K^-_z\},
 \qquad F=I^-_m+b^{h_0}A.
\end{equation}
The relevant dimensions are recorded in
Table~\ref{finite:cancellation-data}.

\begin{table}[htbp]
\centering
\(
\begin{array}{c|rrr}
 m&\dim A&\dim F&\dim(F+K^+_m)\\\hline
 (14,84)&4&5&5\\
 (18,85)&6&5&5
\end{array}
\)
\caption{Dimensions of the cancellation subspaces.}
\label{finite:cancellation-data}
\end{table}

For the second row there is a useful simplification of the product
calculation.  Set $W=I^-_{15,92}+b^3K^-_{9,83}$.  The same column
reductions give
\begin{equation}\label{finite:factorization-containment}
 \dim W=3,\quad
 \rank\bigl(a\colon D_{14,91}\longrightarrow D_{15,92}\bigr)=3,
 \quad \dim\bigl(W+aD_{14,91}\bigr)=3.
\end{equation}
Thus $aD_{14,91}\subseteq W$.  Since $P=aB(t_4')$, its entire
image from $D_{12,76}$ lies in $W$, and $A=K^-_{12,76}$.
The first row uses the product $P\colon D_{4,69}\to D_{7,85}$
directly.  This proves the second row without requiring a separate
matrix for $B(t_4')\colon D_{12,76}\to D_{14,91}$.

Table~\ref{finite:injection-data} gives three tests for injectivity
after taking $L$-homology.

\begin{table}[htbp]
\centering
\(
\begin{array}{ccr|rr}
 f&d&n&\dim I^-_{f,d}&
 \dim\left(I^-_{f,d}+
  \{x\in K^+_{f,d}:b^nx\in I^+_{(f,d)+n|b|}\}\right)\\\hline
 13&71&5&3&3\\
 21&78&3&2&2\\
 14&84&1&4&4
\end{array}
\)
\caption{Enclosure tests for injectivity on fifth-differential homology.}
\label{finite:injection-data}
\end{table}

In each row, an actual cycle whose product by $b^n$ is an actual
boundary lies in the displayed preimage.  Equality of the two
dimensions forces that cycle into $I^-_{f,d}$, so it was already
a boundary.  Hence the three displayed powers of $b$ are injective
on $H$ from their indicated source slots.

Finally, Table~\ref{finite:product-ranks} gives multiplication ranks
on $D$ itself.  The slot $(f,d)$ is the \emph{target}; the source is
$(f-2n,d-3n)$.

\begin{table}[htbp]
\centering
\(
\begin{array}{ccr|rrr}
 f&d&n&\dim\text{ source}&\dim\text{ target}&\rank(b^n)\\\hline
 22&91&2&8&8&8\\
 17&90&2&7&6&6\\
 19&93&1&6&6&6\\
 26&87&6&4&4&4\\
 31&88&6&4&4&4
\end{array}
\)
\caption{Ranks of multiplication by powers of $b$ on $D$.}
\label{finite:product-ranks}
\end{table}

All the displayed subspace calculations use kernels, images,
preimages, and ranks of the specified finite matrices.
Proposition~\ref{alg:table-verification} gives their
construction and reduction over $\F_3$.

Consider $m=(14,84)$ or $(18,85)$ with the data in
\eqref{finite:cancellation-spaces}.  We need two injections on $H$.
First, $b^{h_0}$ is injective from $H_{u+R}$: these are exactly
the first two rows of Table~\ref{finite:injection-data}, with sources
$(13,71)$ and $(21,78)$.  Second, $b^\ell$ is injective from
$H_m$.  For $m=(14,84)$ this is the third row of that table.
For $m=(18,85)$, the first two rows of
Table~\ref{finite:product-ranks} and
Lemma~\ref{finite:homology-transfer}\textup{(ii)} imply that $b^2$
is injective from $H_{18,85}$ into $H_{22,91}$: it is injective
into $D_{22,91}$ and surjective into $D_{17,90}$.
Multiplication by one further $b$ is injective from $H_{22,91}$
into $H_{24,94}$.  Indeed, it is surjective into $D_{19,93}$ by
the third product-rank row and injective into $D_{24,94}$ by
Lemma~\ref{finite:highweight}, since
\[
 94<I_1(24)=98.
\]
Their composite is the required injection by $b^3$ from $H_m$.

Let $c\in H_m$ be a $d_9$-cycle.  Equality of the last two
dimensions in Table~\ref{finite:cancellation-data} gives
$K^+_m\subseteq F$.  Any representative of $c$ is an actual
$L$-cycle, hence belongs to $K^+_m$; its component in $I^-_m$ is
an actual boundary.  Thus
\[
 c=b^{h_0}[x]\quad\text{for some }x\in A\subseteq K^-_u.
\]
By the definition of $A$ there is $w\in K^-_z$ with
\[
 P[x]=b^\ell[w]\quad\text{in }H_p.
\]
Both $x$ and $w$ are actual $L$-cycles because they lie in lower
kernel bounds.  The first required injection now gives
\[
 0=d_9c=b^{h_0}d_9[x]\quad\Longrightarrow\quad d_9[x]=0.
\]
Apply $d_9$ to the preceding product identity.  Using
\eqref{finite:classical-differentials}, permanence of $b$, and the
Leibniz rule yields
\[
 b^\ell d_9[w]=d_9(P[x])=\epsilon'b^6[x]
                         =\epsilon'b^\ell c.
\]
Here $z+R=m$, since $|P|+R=6|b|$.
The second required injection therefore implies
$d_9[w]=\epsilon'c$.  Since $\epsilon'$ is a unit, $c$ is a
boundary.  This proves vanishing at $m=(14,84)$ and $(18,85)$.

At $m=(26,87)$, the last two rows of
Table~\ref{finite:product-ranks} and
Lemma~\ref{finite:homology-transfer}\textup{(i)} show that $b^6$
surjects onto $H_m$: it surjects onto $D_{26,87}$ and injects into
$D_{31,88}$.  To see that $b^6$ is injective on $H$ into
$m+R=(35,89)$, apply
Lemma~\ref{finite:homology-transfer}\textup{(ii)} and the strict
inequalities
\[
 89<I_6(35)=\frac{235}{2},\qquad
 88<S_6(30)=\frac{199}{2}.
\]
Thus a $d_9$-cycle $c\in H_m$ can be written $c=b^6x$; its
differential vanishes only if $d_9x=0$.  Equation
\eqref{finite:classical-differentials} then gives
$d_9(Px)=\epsilon'c$, proving vanishing at $(26,87)$.

It remains to treat $f\geq30$ with $f\equiv2\pmod4$ and
$m=(f,d)$, where $d=(322+f)/4$.  Lemma~\ref{finite:highweight}
gives all four maps needed to repeat the same argument:
\begin{align*}
 b^6\text{ surjects into }D_m
   &\quad\text{if }17f>464,\\
 b^6\text{ injects into }D_{m+J}
   &\quad\text{if }17f>396,\\
 b^6\text{ injects into }D_{m+R}
   &\quad\text{if }17f>328,\\
 b^6\text{ surjects into }D_{m+R-J}
   &\quad\text{if }17f>396.
\end{align*}
Each inequality holds for $f\geq30$.  Hence $b^6$ is surjective
into $H_m$ and injective into $H_{m+R}$ by
Lemma~\ref{finite:homology-transfer}.  For any $d_9$-cycle
$c=b^6x$, injectivity forces $d_9x=0$, and then
$d_9(Px)=\epsilon'c$.  This proves the remaining vanishings and
completes the proof of Lemma~\ref{finite:main}.
